\documentclass[10pt,a4paper]{amsart}
\usepackage[margin=19mm]{geometry}
\usepackage{amsmath,amssymb,mathtools}
\usepackage[scr=rsfs,cal=euler]{mathalfa}
\usepackage{xfrac}
\usepackage[unicode,hidelinks,bookmarks=false]{hyperref}
\newcommand{\ie}{i.e.\ }
\newcommand{\cf}{cf.\ }
\newcommand{\resp}{respectively\ }
\newcommand{\Subsection}{\subsection}
\providecommand{\nobreakdash}{\nobreak\hbox{-}}
\setlength{\emergencystretch}{2em}
\multlinegap0pt
\usepackage{xcolor}
\usepackage{amssymb}
\usepackage{algorithm}\usepackage{algpseudocode}
\usepackage{dsfont}
\usepackage{bm}
\newtheorem{lemma}{Lemma}
\newcommand{\R}{\mathbb{R}} 
\newcommand{\inner}[2]{\left\langle #1,#2\right\rangle}
\newcommand{\one}{\mathbf{1}}
\newcommand{\softmin}{\operatorname{smin}}
\title{Log-Averaged Mirror Prox for Fast, Large-Scale Optimal Transport in Linear Space}
\author{Matthew X. Burns, Jiaming Liang}
\date{Source: arXiv:2511.11359v4 (CC BY 4.0)}
\begin{document}
\maketitle

\noindent\emph{Source and licence.} Log-Averaged Mirror Prox for Fast, Large-Scale Optimal Transport in Linear Space. Version arXiv:2511.11359v4 (CC BY 4.0), distributed by arXiv under the Creative Commons Attribution 4.0 International licence, http://creativecommons.org/licenses/by/4.0/, as recorded in the arXiv licence metadata for this exact version and shown on the abstract page https://arxiv.org/abs/2511.11359v4. Full licence notice: \texttt{licenses/arxiv-2511.11359-CC-BY-4.0.txt}.\par\smallskip

\noindent\textit{Editorial scope.} Counted unit: the article's lemma lem:dual\_infty\_bound (source0.tex lines 800--806). The article's complete original proof of it is delivered with it (source0.tex lines 807--860, one \texttt{proof} environment of 54 lines). No mathematics is written, repaired or re-derived: the statement and the proof are the article's own lines, in the article's own notation, and no retained line is substituted or re-flowed.  Retained uncounted as the source-local context the proof needs: lines 662--664; lines 750--756.  Nothing was omitted from any retained span, so the package carries no omission record.  No retained line contains a citation, so the wrapper carries no bibliography and no bibliography entry.  No retained line contains a figure environment, an image-inclusion command or drawing code, so no image is delivered and the package declares no asset dependency.  Editorial formatting only, in addition: an AMS \texttt{amsart} preamble that restates the article's own declarations and macros used by the retained lines, namely the environments \texttt{lemma} on separate counters, together with the standard packages those lines need.  The retained environments print in this document as Lemma 1 in order of appearance, whereas the article prints the same items with the numbering of its own full document.  The complete native source of the article as distributed in the arXiv e-print for this exact version is bundled unchanged as \texttt{sources/189120/source0.tex}.
\begin{equation}\label{def:dual_eot}
    \max_{\varphi\in\R^n,\psi\in\R^m} \left\{d^{\eta }(\varphi,\psi):=\softmin^{\eta }_{ij}\{C_{ij}+\varphi_i+\psi_j\}-\inner{r}{\varphi} - \inner{c}{\psi}\right\},
\end{equation}

\begin{lemma}\label{lem:logsumexp_additive}
    Let $\varphi\in\R^n$, $\psi\in \R^m $ be two dual potentials. Then, for any $a\in\R$, $b\in\R$,
    \[
    d^{\eta }(\varphi + a\one_n,\psi+b\one_m)=d^{\eta }(\varphi,\psi),
    \]
    where $d^{\eta }$ is the EOT dual defined in~\eqref{def:dual_eot}.
\end{lemma}

\begin{lemma}\label{lem:dual_infty_bound}
    For the dual EOT problem defined in~\eqref{def:dual_eot}, there exists an optimal solution $(\varphi^*,\psi^*)$ such that
    \begin{align*}
        \|\varphi^*\|_\infty&\leq \|C\|_\infty -\eta \log\underset{1\leq i\leq n}\min\{r_i\},\\
        \;\|\psi^*\|_\infty&\leq \|C\|_\infty -\eta \log\underset{1\leq j\leq m}\min\{c_j\}.
    \end{align*}
\end{lemma}

\begin{proof}
    We first claim that there exist dual variables $\varphi^*$ and $\psi^*$ which are optimal solutions to~\eqref{def:dual_eot} satisfying
    \begin{align}
        \nonumber\min_i\varphi^*_i\leq 0\leq \max_i\varphi^*_i,\\
        \min_j\psi^*_j\leq 0\leq \max_j\psi^*_j.\label{ineq:comp_bounds}
    \end{align}
    Let $\hat\varphi^*$, $\hat\psi^*$ be a pair of optimal dual potentials. Define \[
        \Delta_\varphi=\frac{1}{2}(\max_i\hat\varphi_i^*+\min_i\hat\varphi_i^*),\quad
        \Delta_\psi=\frac{1}{2}(\max_j\hat\psi_j^*+\min_j\hat\psi_j^*),
    \]
    and set $\varphi^*=\hat\varphi^* - \Delta_\varphi\one_n$ and $\psi^*=\hat\psi^* - \Delta_\psi\one_m$.
    By Lemma~\ref{lem:logsumexp_additive}, $\varphi^*$ and $\psi^*$ are optimal dual potentials. Furthermore, by construction, they satisfy~\eqref{ineq:comp_bounds}.
    Now, taking the gradient of the dual objective in~\eqref{def:dual_eot}, we have
    \begin{equation}\label{eq:grad_dual_opt}
        0=-r_i +\exp[-\eta ^{-1}\varphi^*_i]\sum_{j=1}^m\frac{\exp[-\eta ^{-1}(C_{ij}+\psi^*_j)]}{Z},
    \end{equation}
    where
    \[Z=\sum_{k,\ell}\exp[-\eta ^{-1}(C_{k\ell}+\varphi^*_k+\psi^*_\ell)].\]
    Rearranging~\eqref{eq:grad_dual_opt}, we obtain
    \begin{equation*}
        \varphi^*_i=-\eta \log r_i + \eta \log\left(\sum_{j=1}^m\exp[-\eta ^{-1}(C_{ij}+\psi^*_j)]\right) - \eta \log Z.
    \end{equation*}
    Since each $C_{ij}\in[0, \|C\|_\infty]$ we have for all $i$, $j$
    \begin{equation*}
    \exp[-\eta ^{-1}(C_{ij}+\psi^*_j)]\geq \exp[-\eta ^{-1}(\psi^*_j)]\exp[-\eta ^{-1}\|C\|_\infty].
    \end{equation*}
    We also have $-\log r_i\geq 0$, hence
    \begin{equation*}
        \varphi^*_i\geq  \eta \log\left(\sum_{j=1}^m\exp[-\eta ^{-1}\psi^*_j]\right) -\|C\|_\infty - \eta \log Z.
    \end{equation*}
    Again using $C_{ij}\geq 0$, we have
    \begin{equation*}
        \varphi^*_i\leq -\eta \log r_i + \eta \log\left(\sum_{j=1}^m\exp[-\eta ^{-1}\psi^*_j]\right) - \eta \log Z.
    \end{equation*}
    Then we have
    \begin{align*}
        \max_{i}\varphi_i^*&\leq -\eta \min_i\log r_i+\eta \log\left(\sum_{j=1}^m\exp[-\eta ^{-1}\psi^*_j]\right)-\eta \log Z\\
        \min_{i}\varphi_i^*&\geq  \eta \log\left(\sum_{j=1}^m\exp[-\eta ^{-1}\psi^*_j]\right) - \|C\|_\infty - \eta \log Z.
    \end{align*}
    So,
    \begin{align*}
        \max_{i}\varphi_i^*-\min_{i}\varphi_i^*\leq \|C\|_\infty-\eta \log\min_{i} r_i.
    \end{align*}
    Since $\min_{i}\varphi_i^*\leq 0$ and $\max_{i}\varphi_i^*\geq 0$, we have that 
    \begin{align*}
        0\geq\min_{i}\varphi_i^*\geq -(\|C\|_\infty-\eta \log\min_{i} r_i),\\
        0\leq\max_{i}\varphi_i^*\leq \|C\|_\infty-\eta \log\min_{i} r_i,
    \end{align*}
    which implies that 
    \[
    \|\varphi^*\|_\infty \leq \|C\|_\infty-\eta \log\min_{i} r_i
    \]
    as claimed. A similar argument holds for the dual variable $\psi^*$ with the $c$ marginal in place of the $r$ marginal.  
\end{proof}

\end{document}
